\documentclass[a4paper,11pt]{article}
% Shared preamble — % Shared preamble — % Shared preamble — \input{../shared/preamble} in every document
% Each document must declare \documentclass[a4paper,11pt]{article} before \input-ing this file.
\usepackage[T1]{fontenc}
\usepackage[utf8]{inputenc}
\usepackage[polish,english]{babel}
\usepackage[top=2cm,bottom=2cm,left=2cm,right=2cm]{geometry}
\usepackage{booktabs}
\usepackage{array}
\usepackage{tabularx}
\usepackage{longtable}
\usepackage{multirow}
\usepackage{hhline}
\usepackage{enumitem}
\usepackage{xcolor}
\usepackage{fancyhdr}
\usepackage{graphicx}
\usepackage{lastpage}

\definecolor{pzzblue}{RGB}{0,70,127}

\pagestyle{fancy}
\fancyhf{}
\renewcommand{\headrulewidth}{0.4pt}
\fancyfoot[C]{\thepage/\pageref{LastPage}}

% Helper: underlined fill-in field of given width
\newcommand{\field}[1]{\underline{\hspace{#1}}}
\newcommand{\fieldline}{\noindent\rule{\linewidth}{0.4pt}}
 in every document
% Each document must declare \documentclass[a4paper,11pt]{article} before \input-ing this file.
\usepackage[T1]{fontenc}
\usepackage[utf8]{inputenc}
\usepackage[polish,english]{babel}
\usepackage[top=2cm,bottom=2cm,left=2cm,right=2cm]{geometry}
\usepackage{booktabs}
\usepackage{array}
\usepackage{tabularx}
\usepackage{longtable}
\usepackage{multirow}
\usepackage{hhline}
\usepackage{enumitem}
\usepackage{xcolor}
\usepackage{fancyhdr}
\usepackage{graphicx}
\usepackage{lastpage}

\definecolor{pzzblue}{RGB}{0,70,127}

\pagestyle{fancy}
\fancyhf{}
\renewcommand{\headrulewidth}{0.4pt}
\fancyfoot[C]{\thepage/\pageref{LastPage}}

% Helper: underlined fill-in field of given width
\newcommand{\field}[1]{\underline{\hspace{#1}}}
\newcommand{\fieldline}{\noindent\rule{\linewidth}{0.4pt}}
 in every document
% Each document must declare \documentclass[a4paper,11pt]{article} before \input-ing this file.
\usepackage[T1]{fontenc}
\usepackage[utf8]{inputenc}
\usepackage[polish,english]{babel}
\usepackage[top=2cm,bottom=2cm,left=2cm,right=2cm]{geometry}
\usepackage{booktabs}
\usepackage{array}
\usepackage{tabularx}
\usepackage{longtable}
\usepackage{multirow}
\usepackage{hhline}
\usepackage{enumitem}
\usepackage{xcolor}
\usepackage{fancyhdr}
\usepackage{graphicx}
\usepackage{lastpage}

\definecolor{pzzblue}{RGB}{0,70,127}

\pagestyle{fancy}
\fancyhf{}
\renewcommand{\headrulewidth}{0.4pt}
\fancyfoot[C]{\thepage/\pageref{LastPage}}

% Helper: underlined fill-in field of given width
\newcommand{\field}[1]{\underline{\hspace{#1}}}
\newcommand{\fieldline}{\noindent\rule{\linewidth}{0.4pt}}


\fancyhead[L]{\textbf{ODPRAWA BEZPIECZEŃSTWA} — Śródlądzie}
\fancyhead[R]{\today}

\begin{document}
\selectlanguage{polish}

\begin{center}
    {\Large\bfseries\color{pzzblue} ODPRAWA BEZPIECZEŃSTWA}\\
    {\large SAFETY BRIEFING}\\[4pt]
    {\normalsize Wody śródlądowe \textbullet{} jacht: \field{5cm} \textbullet{} data: \field{3cm}}
\end{center}

\vspace{0.3cm}
\noindent\rule{\linewidth}{1pt}

\section*{1. Jacht i wyposażenie}

\begin{itemize}[noitemsep]
    \item Gaśnica: \field{5cm} (lokalizacja)
    \item Koło ratunkowe: \field{4cm}
    \item Kamizelki ratunkowe: szafa \field{3cm} — obowiązek dzieci i osób nieumiejących pływać
    \item Apteczka: \field{4cm}
    \item Pompa zęzowa: \field{4cm}
    \item Latarka: \field{3cm}
    \item Kotwica + kotwicznik: \field{4cm}
\end{itemize}

\section*{2. Procedury alarmowe}

\subsection*{Człowiek za burtą (MOB)}
\begin{enumerate}[noitemsep]
    \item Krzyknąć \textbf{„CZŁOWIEK ZA BURTĄ!"}.
    \item Rzucić koło ratunkowe lub boję.
    \item Wyłączyć silnik — uwaga na śrubę!
    \item Podejść do osoby od strony zawietrznej / z prądem.
    \item Wezwać pomoc: telefon 112 / WOPR \field{5cm}.
\end{enumerate}

\subsection*{Wywrotka / przechył}
\begin{enumerate}[noitemsep]
    \item Nie panikować — trzymaj się jachtu.
    \item Sprawdzić stan załogi.
    \item Sygnalizować pomoc (krzyk, świst, machanie).
    \item Telefon 112 / WOPR.
\end{enumerate}

\section*{3. Komunikacja i alarmowanie}

\begin{itemize}[noitemsep]
    \item Telefon alarmowy: \textbf{112}
    \item WOPR (regionalne): \field{6cm}
    \item Straż Pożarna: 998
    \item Bosmanat / Kapitanat: \field{6cm}
    \item Radio VHF (jeśli dostępne): Ch~\field{2cm}
\end{itemize}

\section*{4. Specyfika wód śródlądowych}

\begin{itemize}[noitemsep]
    \item Ograniczone głębokości — nie forsuj nieznanych szlaków nocą.
    \item Śluzy: wyłącz silnik zgodnie z sygnałem śluzowego; liny trzymaj luźno.
    \item Sieci rybackie i bojki — redukuj prędkość w rejonach połowów.
    \item Szybkie łodzie motorowe — zachowaj odległość bezpieczną od toru motorowodnego.
    \item Nagłe burze letnie — obserwuj niebo, wróć do przystani przy ciemniejącym horyzoncie.
    \item Zakaz cumowania do niektórych brzegów — stosuj się do znaków.
\end{itemize}

\section*{5. Zasady ogólne}

\begin{itemize}[noitemsep]
    \item Dzieci na pokładzie — zawsze w kamizelce.
    \item Nie stać na dziobie przy manewrach wyjścia / wejścia.
    \item Zakaz skakania do wody bez zgody kapitana.
    \item Wyłączyć silnik przed wejściem do śluzy.
    \item Lokalizacja wyłącznika głównego: \field{4cm}
    \item Lokalizacja zaworów dennych: \field{5cm}
\end{itemize}

\vspace{0.6cm}
\noindent\rule{\linewidth}{0.4pt}
\vspace{0.2cm}

\noindent\textbf{Potwierdzenie uczestnictwa w odprawie:}

\vspace{0.3cm}
\begin{tabularx}{\linewidth}{|X|X|c|}
\hline
\textbf{Imię i nazwisko} & \textbf{Data} & \textbf{Podpis} \\
\hline
 & & \\[5pt]\hline
 & & \\[5pt]\hline
 & & \\[5pt]\hline
 & & \\[5pt]\hline
 & & \\[5pt]\hline
 & & \\[5pt]\hline
 & & \\[5pt]\hline
 & & \\[5pt]\hline
\end{tabularx}

\vspace{0.5cm}
\noindent Kapitan: \field{6cm} \hfill Podpis: \field{5cm}

\end{document}
